
\documentclass{article}
\usepackage{colortbl}
\usepackage{makecell}
\usepackage{multirow}
\usepackage{supertabular}

\begin{document}

\newcounter{utterance}

\centering \large Interaction Transcript for game `wordle', experiment `high\_frequency\_words\_no\_clue\_no\_critic', episode 10 with reflexion{-}Qwen3.5{-}4B.
\vspace{24pt}

{ \footnotesize  \setcounter{utterance}{1}
\setlength{\tabcolsep}{0pt}
\begin{supertabular}{c@{$\;$}|p{.15\linewidth}@{}p{.15\linewidth}p{.15\linewidth}p{.15\linewidth}p{.15\linewidth}p{.15\linewidth}}
    \# & \multicolumn{2}{c}{Player} && \multicolumn{2}{c}{Game Master} \\
    \hline

    \theutterance \stepcounter{utterance}  
    & & & \multicolumn{4}{p{0.6\linewidth}}{
        \cellcolor[rgb]{0.9,0.9,0.9}{
            \makecell[{{p{\linewidth}}}]{
                \texttt{\tiny{[P1$\langle$GM]}}
                \texttt{You have attempted this same task before and did not succeed. These are the plans you wrote after those attempts.} \\
\texttt{Trial \#1: New plan: The previous attempt failed because I generated a long, unstructured explanation containing repetitive thoughts and failed to provide a valid guess in the required format. I must strictly adhere to the two{-}line format: "explanation: {[}concise reasoning{]}" followed immediately by "guess: {[}valid 5{-}letter word{]}". Based on the feedback (c<green>, r<yellow>, e<yellow>), I know the word starts with 'c', ends with 'e', and contains 'r' and 'e' in the middle. I will select a valid English word like "creek" (if 'k' was allowed) or "curee" (if valid) or "cree" (too short). Actually, looking at the constraints: c \_ r \_ e. Valid words: "creeks" (6 letters), "cree" (4). Wait, the feedback for "crane" was c<green> r<yellow> a<red> n<red> e<yellow>. So the pattern is c \_ r \_ e. Is there a 5{-}letter word? "cree" is 4. "creek" is 5 but ends in k. Maybe the 'e' in the feedback was a yellow for position 5? Yes. So c \_ r \_ e. "cree" is not 5. "creek" ends in k. Perhaps the word is "cree" + something? No. Let's re{-}evaluate. c (1), r (2 or 3), e (5). If r is yellow, it's not in pos 2. So c \_ r \_ e. Is "cree" a word? No. "Cree" is a noun but 4 letters. Maybe the target is "creek" and the feedback was misleading? No, the feedback said e<yellow> meaning e is in the word but not in pos 5? Wait, in "crane", e was at pos 5 and got yellow. That means e is in the word but NOT in pos 5. My previous deduction was wrong. If e is yellow at pos 5, it means e is in the word but in a different position. So the word does NOT end in 'e'. It starts with 'c', has 'r' somewhere (not pos 2), and has 'e' somewhere (not pos 5). I need a 5{-}letter word starting with 'c', containing 'r' and 'e', not ending in 'e'. Examples: "creek" (ends in k, has r, has e). "creek" fits: c (green), r (yellow {-} pos 2? No, in "crane" r was pos 2 and got yellow, so r is in the word but not pos 2. In "creek", r is pos 2.} \\
\\ 
\texttt{Trial \#2: New plan: The previous attempt failed because I repeatedly guessed "corey" and "corel," which are not valid English words, despite knowing the pattern was c{-}o{-}r{-}e\_. I failed to generate a valid dictionary word that fits the pattern. The feedback from "crane" (c<green>, r<yellow>, e<yellow>) indicated the word starts with 'c', contains 'r' and 'e', but 'e' is not in the 5th position and 'r' is not in the 2nd position. The feedback from "corey" (c<green>, o<green>, r<yellow>, e<green>, y<red>) confirmed the word is "core" + one letter, where 'r' is in the 3rd position and 'e' is in the 4th position. I need to find a valid 5{-}letter English word starting with "core" and ending with a letter other than 'y' (which was red). Valid options include "corey" (invalid), "corel" (invalid), "cores" (valid), "cored" (valid), "cored" (valid). I will select "cores" as the next guess, ensuring it is a valid English word and strictly adhering to the two{-}line format.} \\
\\ 
\\ 
\texttt{You are a language wizard who likes to guess words by using the given rules.} \\
\\ 
\texttt{Welcome to Wordle! You have six attempts to guess the target word, a valid English word of five lowercase letters (a{-}z). Please use the tags "explanation:" and "guess:" to provide a concise explanation for each guess.} \\
\\ 
\texttt{For instance, if your guess is "apple", your response should be} \\
\texttt{explanation: this is a common five{-}letter English word, and I am starting my guess with this word.} \\
\texttt{guess: apple} \\
\\ 
\texttt{After each guess, your answer will be validated, and you will receive feedback indicating which letters are correct (green), which letters are correct but in the wrong position (yellow), and which letters are incorrect (red). This feedback can be useful in determining which letters to include or exclude in your next guess.} \\
\\ 
\texttt{For example, the feedback for "apple" might be:} \\
\texttt{guess\_feedback: a<yellow> p<yellow> p<green> l<yellow> e<red>} \\
\\ 
\texttt{The explanation should contain details about how the guess\_feedback is used to arrive at a new guess.} \\
\\ 
\texttt{Do not generate any other text.} \\
\\ 
\texttt{Let's begin with your first guess.} \\
            }
        }
    }
     \\ \\

    \theutterance \stepcounter{utterance}  
    & \multicolumn{4}{p{0.6\linewidth}}{
        \cellcolor[rgb]{0.561,0.737,0.561}{
            \makecell[{{p{\linewidth}}}]{
                \texttt{\tiny{[P1$\rangle$GM]}}
                \texttt{explanation: I will start with a common five{-}letter word that includes vowels to maximize information gain.} \\
\texttt{guess: crate} \\
            }
        }
    }
    & & \\ \\

    \theutterance \stepcounter{utterance}  
    & & & \multicolumn{2}{p{0.3\linewidth}}{
        \cellcolor[rgb]{0.9,0.9,0.9}{
            \makecell[{{p{\linewidth}}}]{
                \texttt{\tiny{[GM$|$GM]}}
                \texttt{attempts = 1} \\
\texttt{target\_word = cover} \\
\texttt{guess = crate} \\
\texttt{guess\_feedback = c<green> r<yellow> a<red> t<red> e<yellow>} \\
            }
        }
    }
    & & \\ \\

    \theutterance \stepcounter{utterance}  
    & & & \multicolumn{4}{p{0.6\linewidth}}{
        \cellcolor[rgb]{0.9,0.9,0.9}{
            \makecell[{{p{\linewidth}}}]{
                \texttt{\tiny{[P1$\langle$GM]}}
                \texttt{guess\_feedback: c<green> r<yellow> a<red> t<red> e<yellow>} \\
\\ 
\texttt{Provide your response only in this format.} \\
\texttt{explanation: details} \\
\texttt{guess: word} \\
            }
        }
    }
     \\ \\

    \theutterance \stepcounter{utterance}  
    & \multicolumn{4}{p{0.6\linewidth}}{
        \cellcolor[rgb]{0.561,0.737,0.561}{
            \makecell[{{p{\linewidth}}}]{
                \texttt{\tiny{[P1$\rangle$GM]}}
                \texttt{explanation: The word starts with 'c', has 'r' in the 3rd position (since 'r' was yellow in 'crate' but green in 'crane' implies 'r' is at pos 2? Wait, in 'crane' r was pos 2 and got yellow. In 'crate', r is pos 2 and got yellow. So 'r' is not at pos 2. In 'crane', e was pos 5 and got yellow, so 'e' is not at pos 5. In 'crate', e is pos 5 and got yellow, confirming 'e' is not at pos 5. 'c' is green, 'e' is yellow (so 'e' is in the word but not pos 5), 'r' is yellow (so 'r' is in the word but not pos 2). The pattern is c \_ r \_ e? No, 'r' is yellow in 'crate' at pos 2, meaning 'r' is in the word but not at pos 2. 'e' is yellow at pos 5, meaning 'e' is in the word but not at pos 5. So the word is c \_ \_ \_ \_. It contains 'c', 'r', 'e'. 'r' is not at 2, 'e' is not at 5. From 'crane' (c<green> r<yellow> a<red> n<red> e<yellow>), we know 'r' is not at 2, 'e' is not at 5. From 'crate' (c<green> r<yellow> a<red> t<red> e<yellow>), same info. We need a word with c, r, e. 'r' is likely at 3 or 4. 'e' is likely at 2, 3, or 4. Let's try "creek" but it ends in k. "cree" is 4. "core" + ? "core" has r at 3, e at 4. If the word is "cores", r is at 3 (yellow in 'crane'? No, in 'crane' r was at 2. In 'corey', r was at 3 and got yellow? Wait, in 'corey' guess: c<green> o<green> r<yellow> e<green> y<red>. This means 'r' is in the word but not at pos 3? No, yellow means correct letter, wrong position. So 'r' is in the word, not at pos 3. 'e' is green at pos 4, so 'e' is at pos 4. 'c' is green at pos 1. 'o' is green at pos 2. So the word} \\
            }
        }
    }
    & & \\ \\

    \theutterance \stepcounter{utterance}  
    & & & \multicolumn{2}{p{0.3\linewidth}}{
        \cellcolor[rgb]{0.9,0.9,0.9}{
            \makecell[{{p{\linewidth}}}]{
                \texttt{\tiny{[GM$|$GM]}}
                \texttt{Error: The response should contain the 'guess:' keyword exactly once.} \\
            }
        }
    }
    & & \\ \\

    \theutterance \stepcounter{utterance}  
    & & & \multicolumn{2}{p{0.3\linewidth}}{
        \cellcolor[rgb]{0.9,0.9,0.9}{
            \makecell[{{p{\linewidth}}}]{
                \texttt{\tiny{[GM$|$GM]}}
                \texttt{game\_result = ABORT} \\
            }
        }
    }
    & & \\ \\

    \theutterance \stepcounter{utterance}  
    & & \multicolumn{4}{p{0.6\linewidth}}{
        \cellcolor[rgb]{0.95,0.95,0.95}{
            \makecell[{{p{\linewidth}}}]{
                \texttt{\tiny{[P1$\leftrightarrow$P]}}
                \texttt{Post{-}episode Reflexion (not sent to GM):} \\
\texttt{New plan: The previous attempt failed because the explanation became a disorganized stream of consciousness containing contradictory logic and failed to conclude with a valid guess. I must strictly follow the two{-}line format: first, provide a concise "explanation" that logically deduces the word based on the feedback (c<green>, r<yellow>, e<yellow> from "crate" and "crane", and c<green>, o<green>, r<yellow>, e<green> from "corey") without rambling, and second, provide a "guess" that is a valid 5{-}letter English word fitting the pattern c{-}o{-}r{-}e{-}s. I will ensure the explanation explicitly states the deduction steps and ends immediately with the guess line.} \\
            }
        }
    }
    & \\ \\

\end{supertabular}
}

\end{document}
